
\documentclass{article}
\usepackage{colortbl}
\usepackage{makecell}
\usepackage{multirow}
\usepackage{supertabular}

\begin{document}

\newcounter{utterance}

\centering \large Interaction Transcript for game `wordle{-}crazy', experiment `high\_frequency\_words\_no\_clue\_no\_critic', episode 6 with Bentel rockers\_\_qwen3.5{-}4b\_\_Bentel\_iter.
\vspace{24pt}

{ \footnotesize  \setcounter{utterance}{1}
\setlength{\tabcolsep}{0pt}
\begin{supertabular}{c@{$\;$}|p{.15\linewidth}@{}p{.15\linewidth}p{.15\linewidth}p{.15\linewidth}p{.15\linewidth}p{.15\linewidth}}
    \# & \multicolumn{2}{c}{Player} && \multicolumn{2}{c}{Game Master} \\
    \hline

    \theutterance \stepcounter{utterance}  
    & & & \multicolumn{4}{p{0.6\linewidth}}{
        \cellcolor[rgb]{0.9,0.9,0.9}{
            \makecell[{{p{\linewidth}}}]{
                \texttt{\tiny{[P1$\langle$GM]}}
                \texttt{You are a language wizard who likes to guess words by using the given rules.} \\
\\ 
\texttt{Welcome to Wordle! You have six attempts to guess the target word, a valid English word of five lowercase letters (a{-}z). Please use the tags "explanation:" and "guess:" to provide a concise explanation for each guess.} \\
\\ 
\texttt{For instance, if your guess is "apple", your response should be} \\
\texttt{explanation: this is a common five{-}letter English word, and I am starting my guess with this word.} \\
\texttt{guess: apple} \\
\\ 
\texttt{After each guess, your answer will be validated, and you will receive feedback indicating which letters are correct (yellow), which letters are correct but in the wrong position (black), and which letters are incorrect (purple). This feedback can be useful in determining which letters to include or exclude in your next guess.} \\
\\ 
\texttt{For example, the feedback for "apple" might be:} \\
\texttt{guess\_feedback: a<black> p<black> p<yellow> l<black> e<purple>} \\
\\ 
\texttt{The explanation should contain details about how the guess\_feedback is used to arrive at a new guess.} \\
\\ 
\texttt{Do not generate any other text.} \\
\\ 
\texttt{Let's begin with your first guess.} \\
            }
        }
    }
     \\ \\

    \theutterance \stepcounter{utterance}  
    & \multicolumn{4}{p{0.6\linewidth}}{
        \cellcolor[rgb]{0.561,0.737,0.561}{
            \makecell[{{p{\linewidth}}}]{
                \texttt{\tiny{[P1$\rangle$GM]}}
                \texttt{explanation: I will start with a very common five{-}letter word that contains a variety of frequently used letters to maximize the information gained from the first guess.} \\
\texttt{guess: stare} \\
            }
        }
    }
    & & \\ \\

    \theutterance \stepcounter{utterance}  
    & & & \multicolumn{2}{p{0.3\linewidth}}{
        \cellcolor[rgb]{0.9,0.9,0.9}{
            \makecell[{{p{\linewidth}}}]{
                \texttt{\tiny{[GM$|$GM]}}
                \texttt{attempts = 1} \\
\texttt{target\_word = major} \\
\texttt{guess = stare} \\
\texttt{guess\_feedback = s<purple> t<purple> a<black> r<black> e<purple>} \\
            }
        }
    }
    & & \\ \\

    \theutterance \stepcounter{utterance}  
    & & & \multicolumn{4}{p{0.6\linewidth}}{
        \cellcolor[rgb]{0.9,0.9,0.9}{
            \makecell[{{p{\linewidth}}}]{
                \texttt{\tiny{[P1$\langle$GM]}}
                \texttt{guess\_feedback: s<purple> t<purple> a<black> r<black> e<purple>} \\
\\ 
\texttt{Provide your response only in this format.} \\
\texttt{explanation: details} \\
\texttt{guess: word} \\
            }
        }
    }
     \\ \\

    \theutterance \stepcounter{utterance}  
    & \multicolumn{4}{p{0.6\linewidth}}{
        \cellcolor[rgb]{0.561,0.737,0.561}{
            \makecell[{{p{\linewidth}}}]{
                \texttt{\tiny{[P1$\rangle$GM]}}
                \texttt{explanation: The feedback indicates that 'a', 'r', and 't' are incorrect letters (purple), while 's' and 'e' are present in the word but in the wrong positions (black). Since I initially thought 's' and 'e' were in the wrong positions, I need to reconsider their placement or assume they might be correct but I placed them incorrectly. However, the feedback format provided in the prompt example uses <black> for correct position and <yellow> for correct letter wrong position. The user's feedback shows <black> for 'a' and 'r', which contradicts my initial assumption that they were wrong. Let me re{-}evaluate: if 'a' and 'r' are <black>, they are in the correct positions. But wait, the user's feedback for "stare" was s<purple> t<purple> a<black> r<black> e<purple>. This means 'a' is in the 3rd position and 'r' is in the 4th position, which matches my guess. So 'a' and 'r' are correct letters in the correct positions. 's' and 'e' are present but in the wrong positions. 't' is not in the word. I need to replace 's' and 'e' with other letters that are likely to be in the word. Let's try "star" with different vowels or consonants. Actually, let's try "stare" again but with different letters for 's' and 'e'. Wait, I already used "stare". Let me try a different word with 'a' and 'r' in the 3rd and 4th positions. How about "barr"? No, that's four letters. Let's try "barre". Or "barrs". But I need five letters. Let's try "barrel". No, that's six. Let's try "barron". No, that's six. Let's try "barren". That's five. But 'e' is in the word, so "barren" has two 'e's. Let's try "barred". That's five. But 'd' is a new letter. Let's try "barred" and see if it works. Wait, I need to be more systematic. The word has 'a' in the 3rd position and 'r' in the 4th position. 's' and 'e' are in the word but not in the 1st and 5th positions respectively. So the word is \_ \_ a r \_. Let's try "barr" with different first and last letters. How about "barrs"? No, that's five but 's' is already used. Let's try "barrs" again. Wait, 's' is in the word but not in the 1st position. So the first letter is not 's'. Let's try "barrs" but with 'b' in the first position and 's' in the last position. But 's' is already in the word, so it can be in the last position. Let's try "barrs". But I need to check if "barrs" is a valid word. It is, but it's informal. Let's try "barre" as a valid word. Or "baron". But 'n' is a new letter. Let's try "baron" and see. Wait, I should try "barrs" first. But let's try "barre" instead. Or "barren". But 'e' is already in the word, so "barren" has two 'e's. Let's try "barred". That's five letters. Let's go with "barred".} \\
\texttt{guess: barred} \\
            }
        }
    }
    & & \\ \\

    \theutterance \stepcounter{utterance}  
    & & & \multicolumn{2}{p{0.3\linewidth}}{
        \cellcolor[rgb]{0.9,0.9,0.9}{
            \makecell[{{p{\linewidth}}}]{
                \texttt{\tiny{[GM$|$GM]}}
                \texttt{Error: The length of the guessed word is not 5.} \\
            }
        }
    }
    & & \\ \\

    \theutterance \stepcounter{utterance}  
    & & & \multicolumn{2}{p{0.3\linewidth}}{
        \cellcolor[rgb]{0.9,0.9,0.9}{
            \makecell[{{p{\linewidth}}}]{
                \texttt{\tiny{[GM$|$GM]}}
                \texttt{game\_result = ABORT} \\
            }
        }
    }
    & & \\ \\

\end{supertabular}
}

\end{document}
